\section{Summary} \label{sec:optimisation_summary}

This chapter has developed a scalable, inexact projection framework to compute the optimal transport map parameterisation, which is summarised in Algorithms~\ref{alg:inexact_pgd}--\ref{alg:accelerated_dykstra} and complemented by Algorithm~\ref{alg:fastforward} developed in Chapter~\ref{ch:dykstra}. By integrating stochastic gradient evaluation and iterative hard thresholding into an outer projected gradient descent loop, the architecture navigates flat objective landscapes and enforces structural sparsity on the physical parameters. Furthermore, by dynamically scheduling the inner Dykstra iteration budget and implementing feasibility-based early termination criteria, the framework circumvents the computational overhead associated with continuous exact Euclidean projections. 

Together with the theoretical resolution of Dykstra's stalling phenomenon presented in Chapter~\ref{ch:dykstra} and the physical structure of the Knothe-Rosenblatt map derived in Chapter~\ref{ch:optimal_transport}, this computational architecture completes the mathematical foundation of the thesis. The subsequent chapter will demonstrate the efficiency of this unified framework by deploying it on several examples, from addressing the problem of propagating non-linear satellite uncertainty in cislunar space to tracking chaotic systems in geophysical applications.